\documentclass[11pt]{article}
\usepackage{amsmath,amssymb}
\usepackage[margin=1.2in]{geometry}
\usepackage{booktabs}
\usepackage{hyperref}

\title{An Exact-Linear-Step (SSCS-style) Sampler for the Active Model\\
\large Error Analysis vs. Euler--Maruyama}
\author{}
\date{}

\begin{document}
\maketitle

\section{Setup: forward and reverse SDEs}

From \texttt{compute\_mean} in \texttt{model\_cifar10\_sde\_DDPM.py}, the forward
(noising) process for the active model is the linear system
\begin{equation}
dx = (-kx + \eta)\,dt + \sqrt{2T_p}\,dW_x,
\qquad
d\eta = -\frac{\eta}{\tau}\,dt + \frac{1}{\tau}\sqrt{2T_a}\,dW_\eta,
\label{eq:forward}
\end{equation}
with $dW_x \perp dW_\eta$. Writing $z = (x,\eta)^\top$, this is $dz = Az\,dt + B\,dW$ with
\begin{equation}
A = \begin{pmatrix} -k & 1 \\ 0 & -1/\tau \end{pmatrix},
\qquad
Q \equiv BB^\top = \begin{pmatrix} 2T_p & 0 \\ 0 & 2T_a/\tau^2 \end{pmatrix}.
\end{equation}

The current sampler (\texttt{sampling()}, lines 810--839) discretizes the
score-corrected \emph{reverse}-time SDE
\begin{equation}
dx = \Big[(kx - \eta) + 2T_p F_x\Big]dt + \sqrt{2T_p}\,dW_x,
\qquad
d\eta = \Big[\frac{\eta}{\tau} + \frac{2T_a}{\tau^2}F_\eta\Big]dt + \frac{1}{\tau}\sqrt{2T_a}\,dW_\eta,
\label{eq:reverse}
\end{equation}
where $(F_x,F_\eta)$ is the network's score output, using plain
Euler--Maruyama (EM) with fixed step $dt$.

\section{Splitting: linear part $L$ vs.\ score part $N$}

The reverse drift in \eqref{eq:reverse} splits additively into a \emph{linear}
part and a \emph{nonlinear} (score) part:
\begin{equation}
\underbrace{\begin{pmatrix} kx-\eta \\ \eta/\tau \end{pmatrix}}_{\tilde A z,\ \ \tilde A = -A}
\;+\;
\underbrace{\begin{pmatrix} 2T_p F_x \\ (2T_a/\tau^2) F_\eta \end{pmatrix}}_{\text{Force}(z)}.
\end{equation}
The key fact: the $\tilde A z$ part, together with its own noise, is an
\emph{exactly solvable} linear SDE — the same OU structure as the forward
process in \eqref{eq:forward}, just with the sign of the drift flipped
($\tilde A = -A$). We integrate that part in closed form and reserve
Euler only for the genuinely nonlinear score term.

\subsection{Exact solution of $L$: $dz = \tilde A z\,dt + B\,dW$}

Diagonalizing $\tilde A = \begin{pmatrix} k & -1 \\ 0 & 1/\tau\end{pmatrix}$
(eigenvalues $k,\ 1/\tau$, assumed distinct — same non-degeneracy already
assumed by \texttt{compute\_mean}) gives
\begin{equation}
e^{\tilde A u} =
\begin{pmatrix}
e_1(u) & \dfrac{e_2(u)-e_1(u)}{\Delta} \\[6pt]
0 & e_2(u)
\end{pmatrix},
\qquad
e_1(u) = e^{ku},\quad e_2(u) = e^{u/\tau},\quad \Delta = k - \frac{1}{\tau}.
\end{equation}

\paragraph{Mean.} Over a step of length $s$, starting from $(x_0,\eta_0)$:
\begin{equation}
\bar x(s) = e_1(s)\,x_0 + \frac{e_2(s)-e_1(s)}{\Delta}\,\eta_0,
\qquad
\bar\eta(s) = e_2(s)\,\eta_0.
\label{eq:mean-exact}
\end{equation}
(This is algebraically identical to \texttt{compute\_mean}$(x_0,\eta_0,t{=}{-}s)$ —
the reverse linear flow is the forward linear flow run at negative time,
since $\tilde A = -A$.)

\paragraph{Covariance.} The noise covariance is \emph{not} simply the forward
covariance at $t=-s$, because $Q$ (fixed) enters the Lyapunov integral
$\Sigma(s)=\int_0^s e^{\tilde A u}Q\,e^{\tilde A^\top u}\,du$ with the flipped-sign
propagator. Carrying out the integral (with $q_1=2T_p$, $q_2=2T_a/\tau^2$):
\begin{align}
\Sigma_{\eta\eta}(s) &= \frac{T_a}{\tau}\big(e_2(s)^2-1\big),
\label{eq:cov-etaeta}\\[4pt]
\Sigma_{x\eta}(s) &= \frac{q_2}{\Delta}
\left[\frac{\tau}{2}\big(e_2(s)^2-1\big) - \frac{e_1(s)e_2(s)-1}{k+1/\tau}\right],
\label{eq:cov-xeta}\\[4pt]
\Sigma_{xx}(s) &= \frac{q_1\big(e_1(s)^2-1\big)}{2k}
+ \frac{q_2}{\Delta^2}
\left[\frac{\tau}{2}\big(e_2(s)^2-1\big) - \frac{2\big(e_1(s)e_2(s)-1\big)}{k+1/\tau} + \frac{e_1(s)^2-1}{2k}\right].
\label{eq:cov-xx}
\end{align}
\emph{(Sanity check: as $s\to0$, $\Sigma_{\eta\eta}\to q_2 s$, $\Sigma_{xx}\to q_1 s$,
$\Sigma_{x\eta}=O(s^2)$ — matches independent-noise-source diffusion at leading
order, since cross-correlation can only appear via drift coupling, which is
higher order in $s$.)}

Because $T_p,T_a,k,\tau$ are constants (no time-dependent noise schedule), the
$2\times2$ matrix $[\bar x,\bar\eta] = M(s)[x_0,\eta_0]^\top$ and $\Sigma(s)$
depend \emph{only on the step size $s$}, not on absolute time — so for a fixed
step schedule, $M(dt/2)$ and $\Sigma(dt/2)$ can be precomputed once, outside
the sampling loop.

\section{The Strang-split (SSCS-style) update}

For each outer step $dt$, with half-step $s = dt/2$:
\begin{enumerate}
\item \textbf{Analytic half-step} (exact, no network call): sample
$(x',\eta') \sim \mathcal N\!\big(M(s)\,(x,\eta),\ \Sigma(s)\big)$ via
\eqref{eq:mean-exact}--\eqref{eq:cov-xx} (Cholesky of $\Sigma(s)$, per channel).
\item \textbf{Score Euler step} (one network evaluation, at $(x',\eta')$,
time $t_{\rm curr}-dt/2$), deterministic only:
\begin{equation}
x'' = x' + 2T_p F_x\,dt, \qquad
\eta'' = \eta' + \frac{2T_a}{\tau^2}F_\eta\,dt.
\end{equation}
\item \textbf{Analytic half-step again}, same formulas as step 1, applied to
$(x'',\eta'')$, giving $(x_{\rm next},\eta_{\rm next})$.
\end{enumerate}
This uses exactly one network evaluation per outer step — the same cost as
the current Euler--Maruyama loop.

\section{Error analysis: why this beats Euler--Maruyama}

\subsection{Euler--Maruyama: order and the stiffness problem}

EM applied to \eqref{eq:reverse} has the standard orders
$\text{(weak) }=1$, $\text{(strong)}=1/2$: global weak error $O(dt)$, global
strong error $O(\sqrt{dt})$. But the more important issue here is not the
asymptotic order — it is the \emph{size of the constant}, which blows up as
$\tau\to 0$.

Consider just the deterministic $\eta$-update in isolation:
$d\eta/ds = \eta/\tau$. EM approximates the exact one-step propagator
$e^{dt/\tau}$ by the linearization $1+dt/\tau$. The local (one-step)
truncation error is
\begin{equation}
e^{dt/\tau} - \Big(1+\frac{dt}{\tau}\Big)
= \frac{1}{2}\Big(\frac{dt}{\tau}\Big)^2 + O\!\Big(\big(\tfrac{dt}{\tau}\big)^3\Big).
\label{eq:lte-em}
\end{equation}
This is the usual $O(dt^2)$ local truncation error of a first-order method —
\emph{but} the coefficient scales as $1/\tau^2$. For fixed $dt$, shrinking
$\tau$ inflates this error rapidly: once $dt/\tau \gtrsim 1$, the linear
approximation $1+dt/\tau$ is no longer even close to $e^{dt/\tau}$, and the
scheme is in the stiff/unstable regime (EM's stability threshold for a mode
with rate $1/\tau$ requires $dt/\tau \lesssim 2$ just to stay bounded, let
alone accurate). This is a direct, quantitative explanation for the
tau$=0.25$ vs.\ tau$=0.45$ FID gap observed empirically: at fixed
$dt = T/N_{\rm steps}$, the smaller $\tau$ pushes $dt/\tau$ higher, inflating
exactly this term.

\subsection{Exact linear step: removes the stiffness error entirely}

Because \eqref{eq:mean-exact}--\eqref{eq:cov-xx} are the \emph{exact} solution
of the linear part for \emph{any} step size $s$, this error term vanishes
identically — not asymptotically, exactly — regardless of how large $dt/\tau$
is. There is no discretization bias on the linear/OU part of the dynamics at
any step size; the only question is how much of the reverse SDE's drift is
linear vs.\ score-coupled, and the noise covariance is likewise exact (both
mean and full covariance are closed-form, not approximated).

\subsection{Remaining error: only from the score (splitting) term
\emph{(corrected)}}

\textbf{Caveat this section originally missed.} The claim below of full
second-order accuracy is \emph{not correct as the scheme was described} in an
earlier draft, and the fix requires stating things carefully.

True Strang splitting is $e^{h(L+N)} \approx \varphi_L^{h/2}\circ e^{hN}\circ
\varphi_L^{h/2}$, which requires the middle factor to be the \emph{exact} flow
of $N$. Our middle step is instead a single frozen-force Euler update
$z \mapsto z + hN(z)$ — an $O(h)$ (first-order) approximation of $e^{hN}$, not
the exact flow. Expanding both the scheme and the true solution to $O(h^2)$
(with $z=(x,\eta)$, $L(z)=(kx-\eta,\eta/\tau)$,
$N(z) = (2T_pF_x,\ (2T_a/\tau^2)F_\eta)$) shows every cross term between $L$
and $N$ ($L'L$, $L'N$, $N'L$) matches exactly between the two — this is the
real benefit of symmetric placement: it cancels the $L$--$N$ \emph{splitting
commutator} error that a naive sequential (Lie--Trotter) scheme would incur.
What does \emph{not} cancel is the term intrinsic to how $N$ alone is
integrated:
\begin{equation}
u_3 - z(dt) = -\frac{dt^2}{2}\,N'(z_0)N(z_0) + O(dt^3),
\end{equation}
because the true flow has $e^{dtN}(z_0) = z_0 + dt\,N(z_0) +
\frac{dt^2}{2}N'(z_0)N(z_0) + O(dt^3)$ while Euler keeps only the first term.
No amount of symmetric bracketing repairs an error internal to a single
sub-flow's own integrator — only errors arising from \emph{interaction}
between $L$ and $N$.

Concretely, in the original variables,
\begin{equation}
N'(z)N(z) =
\begin{pmatrix}
4T_p^2 F_x\,\partial_xF_x + \dfrac{4T_pT_a}{\tau^2}F_\eta\,\partial_\eta F_x \\[8pt]
\dfrac{4T_pT_a}{\tau^2}F_x\,\partial_xF_\eta + \dfrac{4T_a^2}{\tau^4}F_\eta\,\partial_\eta F_\eta
\end{pmatrix},
\end{equation}
i.e.\ Jacobian--vector products of the network with itself. Two things are
worth noting about this term: (i) it is genuinely present, so the scheme is
\textbf{globally first order, not second order}; but (ii) it carries no
$1/\tau$ divergence at fixed step size the way EM's linear-term error
\eqref{eq:lte-em} does — the worst power of $\tau$ here is $\tau^{-4}$
multiplying a \emph{bounded} network output, not an exponentially/rationally
diverging linear term. So the practically dangerous, $\tau$-dependent
stiffness error is still fully removed; what remains is a smooth,
$\tau$-independent $O(dt^2)$ error from the score network's self-nonlinearity.

\subsection{Summary \emph{(corrected)}}

\begin{center}
\begin{tabular}{lll}
\toprule
& Euler--Maruyama & Exact-linear + symmetric split \\
\midrule
Linear ($x,\eta$ OU) part error & $O(dt^2/\tau^2)$ per step, stiff & $0$ (exact) \\
$L$--$N$ splitting/commutator error & $O(dt^2)$ (present, unsandwiched) & $0$ (cancelled by symmetry) \\
$N$ self-nonlinearity error & lumped into same term & $O(dt^2)$ per step, $\tau$-independent \\
Global order & $1$ & $1$ \emph{(not 2 — corrected)} \\
Behavior as $\tau\to0$ at fixed $dt$ & degrades (stiffness) & unaffected \\
Network evals per step & 1 & 1 (same cost) \\
\bottomrule
\end{tabular}
\end{center}

\subsection{Getting to genuine second order}

Two concrete routes, both worth weighing against the current 2-NFE Heun
PF-ODE sampler already in the codebase:
\begin{itemize}
\item \textbf{RK2/midpoint $N$-step (2 NFEs per step).} Replace the Euler
substep with a Heun-style predictor--corrector for $N$ alone (evaluate the
score at $z'$, take an Euler predictor, evaluate the score again at the
predicted point, average). This recovers the missing $\frac12 N'N$ term to
leading order and restores genuine local error $O(dt^3)$ / global order $2$,
at double the network cost per step.
\item \textbf{Score extrapolation (1 NFE per step).} Approximate the missing
term using previously computed score values instead of an extra evaluation
at the current step — an Adams--Bashforth-style linear multistep exponential
integrator, in the spirit of DPM-Solver \cite{lu2022dpmsolver} and DEIS
\cite{zhang2022deis} for diffusion ODEs (both of which are built on exactly
this exact-linear-part-plus-extrapolated-nonlinear-part idea). Needs a
warm-start for the first step and some extra bookkeeping to cache scores
across steps, but keeps the 1-NFE-per-step cost of the current sampler.
\end{itemize}

\section{Final-step correction: an explicit Tweedie estimate at $t_{\rm eps}$}

Both active samplers above stop at $t_{\rm eps}=10^{-3}$ (PF-ODE) or $t=dt$
(SDE/SSCS) rather than the untrained $t=0$, for the same reason documented in
\texttt{tweedie\_denoising\_note.tex} for the passive model: the score network
never saw $t\approx0$ during training. Stopping early avoids the
out-of-distribution evaluation, but leaves the state sitting on the noisy
marginal $p_{t_{\rm eps}}(x,\eta\mid x_0,\eta_0)$ rather than a clean estimate
of $x_0$ — genuine residual spread, not sampler-injected noise.

\subsection{Exact correction}

Unlike the passive model, $\eta_0$ is also random (drawn from its own
stationary prior, independent of the data $x_0$), so the relevant identity is
a joint (not scalar) Tweedie/Miyasawa estimate. Writing $w=(x_0,\eta_0)^\top$,
$F=\begin{pmatrix}a&b\\0&c\end{pmatrix}$ (the mean-map from \S2.1), and
$\Sigma_{\rm cond}=M(t_{\rm eps})$ (the network's actual denoising-score-matching
target covariance, from \texttt{compute\_covariance}), the general linear-Gaussian
Tweedie identity $\nabla_z\log p_t(z) = -\Sigma_{\rm cond}^{-1}(z-F\bar w)$ gives,
after inverting $F$ and keeping only the $x_0$ component:
\begin{equation}
\hat x_0 = \frac{1}{a}\left[(x_t + M_{11}F_x + M_{12}F_\eta)
- \frac{b}{c}\left(\eta_t + M_{12}F_x + M_{22}F_\eta\right)\right].
\label{eq:tweedie-active-exact}
\end{equation}
This is exact given exact scores, and was verified numerically against direct
multivariate Gaussian conditioning (synthetic Gaussian $x_0$ prior, so a
ground-truth posterior mean is computable independently) — the formula
reproduced the ground truth to full floating-point precision.

\subsection{Explicit form when $\tau \gg t_{\rm eps}$ and $T_p \ll 1$}

Both hold for every configuration used in this codebase ($T_p=10^{-3}$,
$\tau\in\{0.25,0.45\}$, $t_{\rm eps}=10^{-3}$). Expanding
\eqref{eq:tweedie-active-exact} to leading order in $t_{\rm eps}$
($a\approx1$, $c\approx1$, $b\approx t_{\rm eps}$, $M_{12}=O(t_{\rm eps}^2)$,
$M_{11}\approx2T_pt_{\rm eps}$):
\begin{equation}
\hat x_0 \approx x_t - t_{\rm eps}\,\eta_t + 2T_p\,t_{\rm eps}\,F_x + O(t_{\rm eps}^2).
\label{eq:tweedie-active-explicit}
\end{equation}
Checked numerically against \eqref{eq:tweedie-active-exact} at the true
parameter values: relative error $\sim0.1\%$, consistent with the $O(t_{\rm
eps}^2)$ truncation. The $2T_pt_{\rm eps}F_x$ term costs nothing extra to keep
(the network is already evaluated once more at $t_{\rm eps}$ for this
correction, same as the passive model's Tweedie step), so
\eqref{eq:tweedie-active-explicit} — not a further-truncated
$\hat x_0\approx x_t-t_{\rm eps}\eta_t$ — is what's implemented in
\texttt{model\_cifar10\_sde\_DDPM.py}'s \texttt{\_active\_tweedie\_correction},
called from the PF-ODE and SDE branches of \texttt{sampling()} and from
\texttt{sampling\_sscs()}, gated by a \texttt{tweedie=True} flag matching the
passive model's convention.

\section{Connection to CLD-SGM}

This is the same construction as the Symmetric Splitting CLD Sampler (SSCS)
in Dockhorn et al.\ (ICLR 2022): their auxiliary velocity variable $v$ plays
the role of our $\eta$, their damping/friction parameters play the role of
our $k,\tau$, and their \texttt{analytical\_dynamics} half-steps are the exact
OU propagation derived above, sandwiching a \texttt{euler\_score\_dynamics}
step that only touches the network-dependent force term. The one difference
is that CLD's noise schedule $\beta(t)$ is time-dependent, so their
closed-form coefficients involve $\beta_{\rm int}$; our $T_p,T_a,k,\tau$ are
constant, which makes our half-step transition matrix and covariance
step-size-only (precomputable once). \emph{Note:} Dockhorn et al.\ justify
SSCS by the same stiffness-removal argument used here, not by a claim of
formal second-order accuracy for the composed sampler — the order correction
in \S4.3 applies equally to SSCS itself, not just to this adaptation of it.

\begin{thebibliography}{3}
\bibitem{dockhorn2022cld}
T.~Dockhorn, A.~Vahdat, and K.~Kreis,
\textit{Score-Based Generative Modeling with Critically-Damped Langevin Diffusion},
ICLR 2022.
\bibitem{song2021sde}
Y.~Song, J.~Sohl-Dickstein, D.~P.~Kingma, A.~Kumar, S.~Ermon, and B.~Poole,
\textit{Score-Based Generative Modeling through Stochastic Differential Equations},
ICLR 2021.
\bibitem{strang1968}
G.~Strang,
\textit{On the Construction and Comparison of Difference Schemes},
SIAM J.\ Numer.\ Anal.\ 5(3), 1968.
\bibitem{lu2022dpmsolver}
C.~Lu, Y.~Zhou, F.~Bao, J.~Chen, C.~Li, and J.~Zhu,
\textit{DPM-Solver: A Fast ODE Solver for Diffusion Probabilistic Model Sampling in Around 10 Steps},
NeurIPS 2022.
\bibitem{zhang2022deis}
Q.~Zhang and Y.~Chen,
\textit{Fast Sampling of Diffusion Models with Exponential Integrator},
ICLR 2023.
\end{thebibliography}

\end{document}